\SourceLeaf{272}
$AC=BB-D$; quare haec forma proprie aequivalebit formae $F$. Iam per formam $Axx+2Bxy+Cyy$ repraesentatur numerus $aa'$, faciendo $x=\mu$, $y=0$, quorum valorum divisor comm. max. est $\mu$; quare $aa'$ etiam per formam $F$ repraesentari poterit ita ut valores indeterminatarum habeant divisorem communem maximum $\mu$ (art.~166). Quoties igitur evadit $\mu=1$, $aa'$ per formam $F$ repraesentari poterit tribuendo indeterminatis valores inter se primos, repraesentatioque haec pertinebit ad valorem $B$ expr.\ $\surd D$ (mod.\ $aa'$), ipsis $b$, $b'$ secundum modulos $a$, $a'$ resp.\ congruum. Conditio $\mu=1$ semper locum habet, quando $a$, $a'$ inter se primi sunt; generaliter autem, quando div.\ comm.\ max.\ ipsorum $a$, $a'$ ad $b+b'$ est primus.

\Needspace{11\baselineskip}
\begin{center}
\textit{Compositio ordinum.}
\end{center}

\Article{245}
\textsc{Theorema.} \textit{Si forma $f$ ad eundem ordinem referenda est ut $g$, similiterque $f'$ est ex eodem ordine ut $g'$: forma $F$ ex $f$, $f'$ composita eundem determinantem habebit ex eodemque ordine erit ut forma $G$ ex $g$, $g'$ composita.}

\textit{Dem.} Sint formae $f,f',F=(a,b,c)$, $(a',b',c')$, $(A,B,C)$ resp., ipsarumque determinantes $=d,d',D$. Porro sit numerorum $a,2b,c$ div.\ comm.\ max.\ $=m$; numerorum $a,b,c$ div.\ comm.\ max.\ $=\mathfrak{m}$; similesque significationes habeant $m'$, $\mathfrak{m}'$ respectu formae $f'$, et $M$, $\mathfrak{M}$ respectu formae $F$. Tunc ordo formae $f$ determinabitur per numeros $d,m,\mathfrak{m}$, unde iidem numeri etiam pro forma $g$ valebunt; eadem ratione numeri $d',m',\mathfrak{m}'$ idem erunt pro forma $g'$ quod sunt pro forma $f'$. Iam per art.~235, numeri $D,M,\mathfrak{M}$ determinati sunt per $d,d',m,m',\mathfrak{m},\mathfrak{m}'$; scilicet erit $D$ divisor communis maximus ipsorum $dm'm'$, $d'mm$; $M=mm'$; atque $\mathfrak{M}=\mathfrak{m}\mathfrak{m}'$ (si simul $m=\mathfrak{m}$, $m'=\mathfrak{m}'$) vel $=2\mathfrak{m}\mathfrak{m}'$ (si $m=2\mathfrak{m}$, aut $m'=2\mathfrak{m}'$). Quae proprietates ipsorum $D,M,\mathfrak{M}$, quum inde sequantur, quod $F$ ex $f,f'$ composita est: facile perspicitur, $D,M$ et $\mathfrak{M}$ etiam pro forma $G$ valere, adeoque $G$ esse ex eodem ordine ut $F$. \textit{Q.~E.~D.}

Ex hac ratione ordinem in quo est forma $F$, compositum dicemus ex ordinibus in quibus sunt formae $f,f'$. Ita \textit{e.~g.}\ ex duobus ordinibus proprie primitivis semper compositus est similis ordo; ex proprie primitivo et improprie primitivo, improprie primitivus. --- Simili modo intelligendum est, si ordo aliquis ex pluribus aliis ordinibus compositus vocabitur.
\Needspace{11\baselineskip}
\SourceLeaf{273}
\begin{center}
\textit{Compositio generum.}
\end{center}

\Article{246}
\textsc{Problema.} \textit{Propositis duabus formis primitivis quibuscunque $f,f'$, ex quarum compositione oritur $F$: ex generibus ad quae pertinent $f,f'$ definire genus ad quod referenda erit $F$.}

\textit{Sol.} I. Consideremus primo eum casum ubi ad minimum una formarum $f,f'$ e.~g.\ prior est proprie primitiva, designemusque determinantes formarum $f,f',F$ per $d,d',D$. Tunc $D$ erit divisor communis maximus numerorum $dm'm'$, $d'$, ubi $m'$ est aut $1$ aut $2$, prout forma $f'$ est proprie aut improprie primitiva; $F$ autem in casu illo pertinebit ad ordinem proprie primitivum, in hoc ad improprie primitivum. Iam genus formae $F$ definietur per ipsius characteres particulares, nempe tum respectu singulorum divisorum primorum imparium ipsius $D$, tum, pro quibusdam casibus, respectu numerorum $4$ aut $8$. Hos igitur singulos determinare oportebit.

$1^{\circ}$. Si $p$ est divisor quicunque primus impar ipsius $D$, necessario etiam ipsos $d,d'$ metietur, adeoque etiam inter characteres formarum $f,f'$ occurrent ipsarum relationes ad $p$. Iam si per $f$ repraesentari potest numerus $a$, per $f'$ numerus $a'$: productum $aa'$ repraesentari poterit per $F$. Si itaque tum per $f$, tum per $f'$ repraesentari possunt residua quadratica ipsius $p$ (per $p$ non divisibilia), etiam per $F$ residua quadratica ipsius $p$ repraesentari poterunt, \textit{i.~e.}\ si utraque $f,f'$ habet characterem $Rp$, forma $F$ eundem characterem habebit. Simili ratione $F$ habebit characterem $Rp$, si utraque $f,f'$ habet characterem $Np$; contra $F$ habebit char.\ $Np$, si altera formarum $f,f'$ habet $Rp$, altera $Np$.

$2^{\circ}$. Si in characterem integrum formae $F$ ingreditur relatio ad numerum $4$, talis relatio etiam in characteres formarum $f$, $f'$ ingredi debet. Nam illud tunc tantummodo evenit, quando $D$ est $\equiv0$ aut $\equiv3$ (mod.\ $4$). Quando $D$ per $4$ est divisibilis, etiam $dm'm'$ et $d'$ per $4$ divisibiles erunt, unde statim patet, $f'$ non posse esse improprie primitivam, adeoque esse $m'=1$; hinc tum $d$ tum $d'$ per $4$ divisibiles erunt, et in utriusque characterem ingredietur relatio ad $4$. Quando $D\equiv3$ (mod.\ $4$), metietur $D$ ipsos $d,d'$; quotientes erunt quadrata, adeoque etiam $d,d'$ necessario vel $\equiv0$ vel $\equiv3$ (mod.\ $4$), et inter cha%
\SourceLeaf{274}%
racteres ipsarum $f,f'$ relatio ad $4$. Hinc eodem modo ut in ($1^{\circ}$) sequitur, characterem formae $F$ fore $1,4$, si vel utraque $f,f'$ habeat $1,4$ vel utraque $3,4$; contra characterem formae $F$ fore $3,4$, si altera formarum $f,f'$ habeat $1,4$, altera $3,4$.

$3^{\circ}$. Quando $D$ per $8$ est divisibilis, etiam $d'$ erit; hinc $f'$ certo proprie primitiva, $m'=1$ atque etiam $d$ per $8$ divisibilis; quare inter characteres formae $F$ aliquis e characteribus $1,8$; $3,8$; $5,8$; $7,8$ tunc tantum locum habere potest, si etiam in charactere tum formae $f$, tum formae $f'$ talis relatio ad $8$ adest. Facile autem confirmatur eodem modo ut ante, characterem formae $F$ fore $1,8$, si $f$ et $f'$ respectu ipsius $8$ eundem habeant; characterem formae $F$ fore $3,8$, si altera formarum $f,f'$ habeat $1,8$ altera $3,8$, vel altera $5,8$ altera $7,8$; $F$ habere $5,8$, si $f,f'$ habeant $1,8$ et $5,8$ vel $3,8$ et $7,8$; $F$ habere $.7,8$, si $f$ et $f'$ habeant vel $1,8$ et $7,8$, vel $3,8$ et $5,8$.

$4^{\circ}$. Quando est $D\equiv2$ (mod.\ $8$), erit $d'$ vel $\equiv0$ vel $\equiv2$ (mod.\ $8$), hinc $m'=1$, adeoque etiam $d$ vel $\equiv0$ vel $\equiv2$ (mod.\ $8$); attamen uterque $d,d'$ per $8$ divisibilis esse nequit, quoniam $D$ est divisor communis \textit{maximus} ipsorum. Quare in eo tantum casu alteruter characterum $1\text{ \textit{et} }7,8$; $3\text{ \textit{et} }5,8$, formae $F$ tribui debebit, ubi vel vtraque forma $f,f'$ aliquem ex illis habet, vel altera aliquem ex illis, altera aliquem horum $1,8$; $3,8$; $5,8$; $7,8$. Hinc facile deducitur, characterem formae $F$ determinari per tabulam sequentem, si character in margine positus pertineat ad alteram formarum $f,f'$, ad alteram vero character in facie:
\[
\begin{array}{||c||c|c||}
\hline\hline
 & 1\text{ \textit{et} }7,8 & 3\text{ \textit{et} }5,8 \\
 & \text{vel }1,8 & \text{vel }3,8 \\
 & \text{vel }7,8 & \text{vel }5,8 \\
\hline\hline
1\text{ \textit{et} }7,8 & 1\text{ \textit{et} }7,8 & 3\text{ \textit{et} }5,8 \\
\hline
3\text{ \textit{et} }5,8 & 3\text{ \textit{et} }5,8 & 1\text{ \textit{et} }7.8 \\
\hline\hline
\end{array}
\]

$5^{\circ}$. Eodem modo probatur, ipsi $F$ tribui non posse alterutrum characterum $1\text{ \textit{et} }3,8$; $5\text{ \textit{et} }7,8$, nisi etiam aliquis ex iisdem saltem uni formarum $f,f'$ competat, alterique vel aliquis ex iisdem, vel aliquis ex his $1,8$; $3,8$; $5,8$; $7,8$.
\SourceLeaf{275}
Et quidem character formae $F$ determinabitur per hanc tabulam, in cuius margine et facie sunt characteres formarum $f,f'$
\[
\begin{array}{||c||c|c||}
\hline\hline
 & 1\text{ \textit{et} }3,8 & 5\text{ \textit{et} }7,8 \\
 & \text{vel }1,8 & \text{vel }5,8 \\
 & \text{vel }3,8 & \text{vel }7,8 \\
\hline\hline
1\text{ \textit{et} }3,8 & 1\text{ \textit{et} }3,8 & 5\text{ \textit{et} }7,8 \\
\hline
5\text{ \textit{et} }7,8 & 5\text{ \textit{et} }7,8 & 1\text{ \textit{et} }3.8 \\
\hline\hline
\end{array}
\]

II. Si utraque forma $f,f'$ est improprie primitiva, erit $D$ divisor communis maximus numerorum $4d.\ 4d'$, sive $\frac{1}{4}D$ div.\ comm.\ maximus numerorum $d.\ d'$. Hinc facile sequitur, tum $d$, tum $d'$, tum $\frac{1}{4}D$ fore $\equiv1$ (mod.\ $4$). Ponendo autem $F=(A,B,C)$, div.\ comm.\ max.\ numerorum $A,B,C$ erit $=2$. et div.\ comm.\ max.\ numerorum $A,2B,C$ erit $4$. Quare $F$ erit forma derivata ex improprie primitiva $(\frac{1}{2}A,\frac{1}{2}B,\frac{1}{2}C)$, cuius determinans erit $\frac{1}{4}D$, et cuius genus determinabit genus formae $F$. Character autem illius formae, tamquam improprie primitivae, relationes ad $4$ vel $8$ non implicabit, sed tantummodo relationes ad singulos divisores primos impares ipsius $\frac{1}{4}D$. Iam quum omnes hi divisores manifesto etiam ipsos $d,d'$ metiantur, atque semissis cuiusvis producti duorum factorum. quorum alter per $f$ alter per $f'$ est repraesentabilis, per formam $(\frac{1}{2}A.\ \frac{1}{2}B,\frac{1}{2}C)$ repraesentari possit: facile perspicietur, characterem huius formae respectu cuiusvis numeri primi imparis $p$ ipsum $\frac{1}{4}D$ metientis fore $Rp$, \textit{tum} si fuerit $2Rp$ atque formae $f,f'$ respectu ipsius $p$ eundem characterem habeant. \textit{tum} si fuerit $2Np$ atque characteres formarum $f,f'$ respectu ipsius $p$ oppositi; contra characterem illius formae fore $Np$, \textit{tum} si $f,f'$ habeant characteres aequales respectu ipsius $p$ atque sit $2Np$, \textit{tum} si $f,f'$ habeant oppositos atque sit $2Rp$.

\Needspace{8\baselineskip}
\Article{247}
Ex solutione problematis praec.\ manifestum est, si $g$ sit forma primitiva ex eodem ordine et genere ut $f$, nec non $g'$ forma primitiva ex eodem ordine et genere ut $f'$: formam ex $g$ et $g'$ compositam ad idem genus pertinere, ad quod pertineat forma ex $f$ et $f'$ composita. Hinc sponte sequitur significatio \textit{generis}
\SourceLeaf{276}
ex duobus aliis generibus (sive etiam pluribus) \textit{compositi}. Porro ibinde patet, si $f,f'$ eundem determinantem habeant atque $f$ sit forma e genere principali, $F$ vero ex $f$ et $f'$ composita: $F$ fore ex eodem genere ut $f'$; quocirca genus principale in compositione cum aliis generibus eiusdem determinantis semper omitti poterit. Si vero reliquis manentibus $f$ non est e genere principali, $f'$ autem forma primitiva: $F$ certo erit ex alio genere quam $f'$. Denique si $f,f'$ sunt formae proprie primitivae eiusdem generis, $F$ erit e genere principali; si vero $f,f'$ sunt ambae proprie primitivae eiusdem determinantis, sed e diversis generibus, $F$ ad genus principale pertinere non poterit. Quodsi itaque forma quaecunque proprie primitiva \textit{cum se ipsa} componitur, forma inde resultans, quae etiam proprie primitiva eiusdemque determinantis erit, necessario ad genus principale pertinebit.

\Needspace{8\baselineskip}
\Article{248}
\textsc{Problema.} \textit{Propositis duabus formis quibuscunque $f,f'$, e quibus composita est $F$: e generibus formarum $f,f'$ definire genus formae $F$.}

\textit{Sol.} Sit $f=(a,b,c)$, $f'=(a',b',c')$, $F=(A,B,C)$, porro $\mathfrak{m}$ div.\ comm.\ max.\ numerorum $a,b,c$, atque $\mathfrak{m}'$ div.\ comm.\ max.\ numerorum $a',b',c'$, ita ut $f,f'$ sint derivatae e primitivis $(\frac{a}{\mathfrak{m}},\allowbreak \frac{b}{\mathfrak{m}},\allowbreak \frac{c}{\mathfrak{m}})$, $(\frac{a'}{\mathfrak{m}'},\allowbreak \frac{b'}{\mathfrak{m}'},\allowbreak \frac{c'}{\mathfrak{m}'})$, quas denotabimus per $\mathfrak{f}$, $\mathfrak{f}'$ resp. Iam si saltem una formarum $\mathfrak{f}$, $\mathfrak{f}'$ est proprie primitiva, divisor comm.\ max.\ numerorum $A,B,C$ erit $\mathfrak{m}\mathfrak{m}'$, adeoque $F$ derivata e forma primitiva $(\frac{A}{\mathfrak{m}\mathfrak{m}'},\allowbreak \frac{B}{\mathfrak{m}\mathfrak{m}'},\allowbreak \frac{C}{\mathfrak{m}\mathfrak{m}'})\ldots\mathfrak{F}$, unde patet, genus formae $F$ pendere a genere formae $\mathfrak{F}$. Sed facile perspicietur, $\mathfrak{F}$ per eandem substitutionem transire in $\mathfrak{f}\mathfrak{f}'$, per quam $F$ transeat in $ff'$, adeoque $\mathfrak{F}$ ex $\mathfrak{f}$, $\mathfrak{f}'$ esse compositam, ipsiusque genus per problema art.~246 determinari posse. --- Si vero utraque $\mathfrak{f}$, $\mathfrak{f}'$ est improprie primitiva, divisor c.~m.\ numerorum $A,B,C$ erit $2\mathfrak{m}\mathfrak{m}'$, formaque $\mathfrak{F}$ etiamnum ex $\mathfrak{f}$, $\mathfrak{f}'$ composita et manifesto e proprie primitiva $(\frac{A}{2\mathfrak{m}\mathfrak{m}'},\allowbreak \frac{B}{2\mathfrak{m}\mathfrak{m}'},\allowbreak \frac{C}{2\mathfrak{m}\mathfrak{m}'})$ derivata. Huius itaque formae genus determinari poterit per art.~246; et quum $F$ ex eadem forma derivata sit, ipsius genus hinc sponte innotescit.

Ex hac solutione manifestum est, theorema in art.\ praec.\ pro formis primitivis explicatum, scilicet \textit{si $f',g'$ sint ex iisdem generibus resp.\ ut $f,g$, formam ex $f',g'$ compositam ex eodem genere fore, ex quo sit forma ex $f,g$ composita}, generaliter pro formis quibuscunque valere.
\Needspace{11\baselineskip}
\SourceLeaf{277}
\begin{center}
\textit{Compositio classium.}
\end{center}

\Article{249}
\textsc{Theorema.} \textit{Si formae $f,f'$ sunt ex iisdem ordinibus generibus et classibus ut $g,g'$ resp.: forma ex $f$ et $f'$ composita ex eadem classe erit ut forma ex $g$ et $g'$ composita.}

Ex hoc theoremate (cuius veritas ex art.~239 protinus sequitur) sponte patebit significatio \textit{classis e duabus classibus datis} sive etiam \textit{e pluribus compositae}.

Si classis quaecunque $K$ cum classe principali componitur, classis $K$ ipsa prodibit sive classis principalis in compositione cum aliis classibus eiusdem determinantis negligi potest. Ex compositione duarum classium oppositarum proprie primitivarum semper oritur classis principalis eiusdem determinantis (v.\ art.~243). Quum itaque quaevis classis anceps sibi ipsa opposita sit: ex compositione cuiusvis classis ancipitis proprie primitivae cum se ipsa classis principalis eiusdem determinantis provenit.

Propositio ultima etiam conversa valet: scilicet \textit{si ex compositione classis proprie primitivae $K$ cum se ipsa provenit classis principalis $H$ eiusdem determinantis, $K$ necessario erit classis anceps.} Si enim $K'$ est classis opposita ipsi $K$, e tribus classibus $K,K,K'$ composita erit eadem classis quae oritur ex $H$ et $K'$; ex illis provenit $K$ (quoniam $K$ et $K'$ producunt $H$, haec cum $K$ ipsam $K$), ex his $K'$; quare $K$ cum $K'$ coincidet eritque adeo classis anceps.

Porro notetur propositio haec: \textit{Si classes $K,L$ oppositae sunt classibus $K',L'$ resp.: classis ex $K$ et $L$ composita classi ex $K'$ et $L'$ compositae erit opposita.} Sint formae $f,g,f',g'$ resp.\ e classibus $K,L,K',L'$; forma $F$ composita ex $f,g$, atque $F'$ composita ex $f',g'$. Quum $f'$ ipsi $f$, atque $g'$ ipsi $g$ improprie aequivaleant, $F$ autem composita sit ex utraque $f,g$ directe: $F$ etiam ex $f',g'$ composita erit, sed ex utraque inverse. Quare forma quaecunque, quae ipsi $F$ improprie aequivalet, composita erit ex $f',g'$ directe adeoque ipsi $F'$ proprie aequivalebit (artt.~238.~239), unde $F,F'$ improprie aequivalebunt, classesque ad quas pertinent, oppositae erunt.

Hinc sequitur, si classis anceps $K$ cum classe ancipite $L$ componatur. semper prodire classem ancipitem. Nam opposita erit classi, quae composita est e classibus ipsis $K,L$ oppositis, adeoque sibi ipsi, quoniam hae classes sibi ipsae sunt oppositae.

\SourceLeaf{278}
Denique observamus, si propositae sint classes duae quaecunque $K,L$ eiusdem determinantis, quarum prior sit proprie primitiva, semper inveniri posse classem $M$ eiusdem determinantis, ex qua atque $K$ composita sit $L$. Manifesto hoc obtinetur, accipiendo pro $M$ classem quae composita est ex $L$ atque classe ipsi $K$ opposita; simul perspicietur facillime, hanc classem esse unicam quae hac proprietate sit praedita, sive classes diversas eiusdem det.\ cum eadem classe pr.\ prim.\ compositas producere classes diversas.

Classium compositio commode per signum additionis, $+$, denotari potest, sicuti classium identitas per signum aequalitatis. In his signis propositio modo tradita exhiberi potest ita: Si $K'$ est classis opposita ipsi $K$, erit $K+K'$ classis principalis eiusdem determinantis, unde $K+K'+L=L$; posita itaque $K'+L=M$, erit $K+M=L$, uti desiderabatur; si vero praeter $M$ alia $M$ daretur eadem proprietate praedita, sive $K+M'=L$, foret $K+K'+M'=L+K'=M$, unde $M'=M$. --- Si plures classes identicae componuntur, hoc (ad instar multiplicationis) denotari potest praefigendo ipsarum numerum, ita ut $2K$ idem designet ut $K+K$, $3K$ idem ut $K+K+K$ etc. Eadem signa etiam ad formas transferri possent, ita ut $(a,b,c)+(a',b',c')$ designaret formam ex $(a,b,c)$, $(a',b',c')$ compositam: sed ne vel species ambiguitatis oriri possit, hac abbreviatione abstinere malumus, praesertim quod tali signo $\surd M(a,b,c)$ significationem peculiarem iam tribuimus. --- Classem $2K$ ex \textit{duplicatione} classis $K$ oriri dicemus, classem $3K$ ex \textit{triplicatione} etc.

\Needspace{8\baselineskip}
\Article{250}
Si $D$ est numerus per $mm$ divisibilis (ubi ipsum $m$ positivum supponimus): dabitur ordo formarum determinantis $D$ ex ordine proprie primitivo determinantis $\frac{D}{mm}$ derivatus (sive \textit{duo}, quando $D$ est negativus, nempe positivus et negativus); manifesto forma $(m,0,-\frac{D}{m})$ ad illum ordinem pertinebit (scilicet ad positivum) meritoque tamquam \textit{forma simplicissima} in eo considerari potest (sicuti $(-m,0,\frac{D}{m})$ erit simplicissima in ordine negativo quando $D$ neg.). Si insuper est $\frac{D}{mm}\equiv1$ (mod.\ $4$), dabitur etiam ordo formarum det.\ $D$ ex improprie primitivo det.\ $\frac{D}{mm}$ derivatus, ad quem manifesto forma $(2m,m,\frac{mm-D}{2m})$ pertinebit et pro simplicissima in eodem habebitur. (Quando $D$ est neg., rursus duo ordines dabuntur et in negativo forma $(-2m,-m,\frac{D-mm}{2m})$ pro simplicissima habebitur). Ita \textit{e.~g.}, si etiam eum casum ubi $m=1$ huc referre lubet, in quatuor ordinibus
\SourceLeaf{279}
formarum det.\ $45$ sequentes erunt simplicissimae $(1,0,-45)$, $(2,1,-22)$. $(3,0,-15)$, $(6,3,-6)$. Quibus ita intellectis, offert se

\textsc{Problema.} \textit{Proposita forma quacunque $F$ ex ordine $O$, invenire formam proprie primitivam (positivam) eiusdem determinantis, ex cuius compositione cum forma in $O$ simplicissima oriatur $F$.}

\textit{Sol.} Sit forma $F=(ma,mb,mc)$ derivata e primitiva $f=(a,b,c)$ cuius determinans $=d$, supponamusque \textit{primo}, $f$ esse proprie primitivam. Primo observamus, si forte $a$ ad $2dm$ non sit primus, certo dari alias formas ipsi $(a,b,c)$ proprie aequivalentes, quarum termini primi hac proprietate sint praediti. Nam per art.~228 dantur numeri ad $2dm$ primi per formam illam repraesentabiles; sit talis numerus $a'=a\alpha\alpha+2b\alpha\gamma+c\gamma\gamma$, supponamusque, (quod licet), $\alpha,\gamma$ esse inter se primos; tum, acceptis $\beta,\delta$ ita ut fiat $\alpha\delta-\beta\gamma=1$, transeat $f$ per substitutionem $\alpha,\beta,\gamma,\delta$ in formam $(a',b',c')$, quae illi proprie aequivalebit et proprietate praescripta erit praedita. Iam quum etiam $F$ et $(a'm,b'm,c'm)$ proprie aequivaleant, facile perspicietur, sufficere eum casum considerare ubi $a$ ad $2dm$ sit primus. Tunc $(a,bm,cmm)$ erit forma proprie primitiva (si enim $a,2bm,cmm$ divisorem communem haberent, hunc etiam $2dm=2bbm-2acm$ implicaret) eiusdem determinantis ut $F$, confirmaturque facile, $F$ transmutari in productum e forma $(m,0,-dm)$, quae, nisi $F$ est forma negativa, erit simplicissima ordinis $O$, in $(a,bm,cmm)$ per substitutionem $1,0,-b,-cm;\ 0,m,a,bm$, unde per criterium in obs.~4. art.~235 concluditur, $F$ ex $(m,0,-dm)$ et $(a,bm,cmm)$ esse compositam. Quando autem $F$ est forma negativa, transibit in productum e forma simplicissima eiusdem ordinis $(-m,0,dm)$ in positivam $(-a,bm,-cmm)$ per substitutionem $1,0,b,-cm;\ 0,-m,-a,bm$, adeoque ex ipsis erit composita.

\textit{Secundo}, si $f$ est forma improprie primitiva, supponere licebit $\frac{1}{2}a$ ad $2dm$ esse primum; si enim haec proprietas in forma $f$ locum nondum habet, inveniri potest forma ipsi $f$ proprie aequivalens et hac proprietate praedita. Hinc autem sequitur facile, $(\frac{1}{2}a,bm,2cmm)$ esse formam proprie primitivam eiusdem determinantis ut $F$; aeque facile confirmatur, $F$ transire in productum e formis
\[
(\pm2m,\pm m,\pm\tfrac{1}{2}(m-dm)),\quad
(\pm\tfrac{1}{2}a,bm,\pm2cmm)
\]
per substitutionem
\SourceLeaf{280}
\[
1,0,\tfrac{1}{2}(1\mp b),-cm;\quad 0,\pm2m,\pm\tfrac{1}{2}a,(b\pm1)m
\]
ubi signa inferiora accipienda sunt quando $F$ est forma negativa, superiora in casibus reliquis, adeoque ex his duabus formis esse compositam, quarum prior erit simplicissima ordinis $O$, posterior forma proprie primitiva (positiva).

\Needspace{8\baselineskip}
\Article{251}
\textsc{Problema.} \textit{Propositis duabus formis $F,f$ eiusdem determinantis $D$ et ad eundem ordinem $O$ pertinentibus: invenire formam proprie primitivam determinantis $D$, quae cum $f$ composita producat $F$.}

\textit{Sol.} Sit $\varphi$ forma simplicissima ordinis $O$; $\mathfrak{F}$, $\mathfrak{f}$ formae proprie primitivae det.\ $D$, quae cum $\varphi$ compositae producant ipsas $F.\ f$ resp.; denique $f'$ forma proprie primitiva, quae cum $\mathfrak{f}$ composita producat $\mathfrak{F}$. Tunc forma $F$ composita erit e tribus formis $\varphi,\mathfrak{f},f'$, sive e duabus $f,f'$. \textit{Q.~E.~I.}

Quaevis itaque classis ordinis dati considerari potest tamquam composita ex quacunque classe data eiusdem ordinis et aliqua classe proprie primitiva eiusdem determinantis.

\Needspace{11\baselineskip}
\begin{center}
\textit{Pro determinante dato in singulis generibus eiusdem ordinis contentae sunt classes aeque multae.}
\end{center}

\Article{252}
\textsc{Theorema.} \textit{Pro determinante dato in singulis generibus eiusdem ordinis contentae sunt classes aeque multae.}

\textit{Dem.} Pertineant genera $G$ et $H$ ad eundem ordinem, constet $G$ ex $n$ classibus $K,K',K''\ldots K^{n-1}$, sitque $L$ classis aliqua e genere $H$. Investigetur per art.\ praec.\ classis proprie primitiva $M$ eiusdem determinantis, ex cuius compositione cum $K$ prodeat $L$, designenturque classes quae oriuntur ex compositione classis $M$ cum $K',K''\ldots K^{n-1}$ resp.\ per $L',L''\ldots L^{n-1}$. Tunc ex obs.\ ultima art.~249 sequitur, omnes classes $L,L',L''\ldots L^{n-1}$ esse diversas, et per art.~248 omnes pertinebunt ad genus idem, \textit{i.~e.}\ ad genus $H$. Denique perspicietur facile, $H$ alias classes praeter has continere non posse, quum quaevis classis generis $H$ tamquam composita considerari possit ex $M$ et alia classe eiusdem determinantis, quae necessario semper erit e genere $G$. Quocirca $H$ perinde ut $G$ continet $n$ classes diversas. \textit{Q.~E.~D.}
\Needspace{11\baselineskip}
\SourceLeaf{281}
\begin{center}
\textit{Comparantur multitudines classium in singulis generibus ordinum diversorum contentarum.}
\end{center}

\Article{253}
Theorema praecedens supponit ordinis identitatem neque ad ordines diversos est extendendum. Ita \textit{e.~g.}\ pro determinante $-171$ dantur $20$ classes positivae, quae reducuntur ad quatuor ordines: in ordine proprie primitivo duo continentur genera, utrumque sex classes complectitur; in ordine impr.\ primitivo duo genera quatuor classes possident, singula binas; in ordine derivato ex O.\ proprie prim.\ det.\ $-19$ unicum est genus tres classes complectens; denique O.\ derivatus ex impr.\ prim.\ det.\ $-19$ unicum genus habet ex una classe constans; perinde se habent classes negativae. Operae itaque pretium est, in principium generale inquirere, a quo nexus inter multitudines classium in diversis ordinibus pendeat. Supponamus, $K,L$ esse duas classes ex eodem ordine (positivo) $O$ determinantis $D$, atque $M$ classem proprie primitivam eiusdem det., ex cuius compositione cum $K$ oriatur $L$, qualis per art.~251 semper potest assignari. Iam in quibusdam casibus fieri potest, ut $M$ sit \textit{unica} classis pr.\ primitiva, quae cum $K$ composita producat $L$; in aliis plures classes diversae pr.\ primitivae exstare possunt hac proprietate praeditae. Supponamus generaliter, dari $r$ huiusmodi classes pr.\ primitivas, $M,M',M''\ldots M^{r-1}$, quae singulae cum $K$ compositae producant eandem classem $L$, designemusque illarum complexum per $W$. Porro sit $L'$ alia classis ordinis $O$ (a classe $L$ diversa), atque $N'$ classis pr.\ prim.\ det.\ $D$, quae cum $L$ composita efficiat $L'$, designeturque complexus classium $N'+M,N'+M',N'+M''\ldots N'+M^{r-1}$ (quae omnes erunt proprie primitivae et inter se diversae) per $W'$. Tunc perspicietur facile, $K$ cum classe quacunque ex $W'$ compositam producere $L'$, unde concluditur, $W$ et $W'$ nullam classem communem habere; praeterea nullo negotio comprobatur, nullam classem pr.\ primitivam in complexu $W'$ non contentam dari, quae cum $K$ composita producat ipsam $L'$. Eodem modo patet, si $L''$ sit alia classis ordinis $O$ a classibus $L,L'$ diversa, dari $r$ formas pr.\ primitivas tum inter se tum a formis $W,W'$ diversas, quae singulae cum $K$ compositae ipsam $L''$ producant, et perinde res se habebit pro omnibus reliquis classibus ordinis $O$. Quoniam vero quaevis classis pr.\ prim.\ (positiva) determinantis $D$ cum $K$ composita classem ordinis $O$ producit, facile hinc colligitur, si multitudo omnium classium ordinis $O$ sit $n$, multitudinem omnium classium proprie primitivarum (positivarum) eiusdem determinantis fore $rn$. 
